
\documentclass{article}
\usepackage{hyperref}
\begin{document}

\title{Software Design Document (SDD)}
\author{Team 13}
\date{\today}
\maketitle

\newpage

\tableofcontents

\newpage

\section{Version Description}
Version 1.0 - Initial document setup for Snapshot 1.

\section{Introduction}

\subsection{Purpose}
The purpose of this document is to describe the overall software design and structure of the project. This document provides an overview of the system architecture, user interface, workflow, and project components.

\subsection{Intended Audience}
This document is intended for developers, instructors, testers, and users who want to understand the design and functionality of the application.

\subsection{System Overview}
This project is a mock Final Fantasy VII-style RPG game application. The game allows users to explore a small game environment, interact with characters, and participate in turn-based battles.

\section{System Architecture}

\subsection{Workflow}
The user launches the game application and interacts with the user interface. The game engine processes user input and updates gameplay systems such as movement, combat, and inventory management.

\subsection{Components}
The project is divided into multiple components including:
\begin{itemize}
    \item User Interface
    \item Game Logic
    \item Battle System
    \item Inventory System
    \item Save Data Management
\end{itemize}

\section{User Interface}

\subsection{Usage}
Users interact with the game through keyboard controls and graphical menus. The interface allows users to navigate menus, move characters, and engage in battles.

\subsection{Database Design}
The project may use file-based save storage to maintain player progress, inventory information, and character statistics.

\section{Glossary}

\begin{itemize}
    \item UI - User Interface
    \item RPG - Role Playing Game
    \item DB - Database
\end{itemize}


\section{References}

\begin{itemize}
    \item \href{https://docs.github.com/}{GitHub Documentation}
    \item \href{https://support.atlassian.com/jira-software-cloud/}{Jira Documentation}
    \item \href{https://support.testrail.com/}{TestRail Documentation}
    \item \href{https://www.latex-project.org/help/documentation/}{LaTeX Documentation}
\end{itemize}



\end{document}

